% HandNet-Bench. Owner: claude-13.
\section{HandNet-Bench: the protocol as a benchmark}
\label{sec:benchmark}
The protocol of Section~\ref{sec:protocol} does not depend on who the author is.
Fixing the prompt turns it into a benchmark for language models whose score is the
test accuracy of the network the model writes---a measurement of whether the model's
knowledge of a domain is articulate enough to be written down as a mechanism, which
a question-answering benchmark cannot give.

Every model receives the identical text: the rules (no training, no data, a tensor
forward pass with every behaviour coming from a number in \texttt{WEIGHTS}, every
weight commented, at most $10^5$ parameters) and a task description. It returns one
self-contained Python file. The harness checks the rules statically (forbidden
imports, data reads, references to data files), counts parameters and non-zero
entries, scores the file once on the sealed test set, and records the row. A
bounded revision round gives every model the same development report (for MNIST,
the confusion matrix on 500 labelled training images with class means and two
errors per class as ASCII; for SST-2, accuracy and forty misclassified development
sentences) and scores the revised file as its own row with the labels it saw. The
tasks at release are MNIST, SST-2 and ten-digit addition as a sequence model; a
task is a data loader with a sealed split, a prompt and, optionally, a report
template.

Table~\ref{tab:bench} is the leaderboard at submission. The reference rows are the
networks of this paper (adapters around the repository's modules, marked as not
self-contained). The cross-model rows are the first run of the benchmark proper:
four Claude models---Haiku 5.5, Sonnet 5.5, Opus 5.5 and Fable 5.1---each received
the byte-identical prompt in a fresh context with no tools (the agent harness
reports zero tool calls for every row) and returned one self-contained file,
which was scored once. Every model chose a template matcher on a $2\times2$-pooled
image; they differ in how the templates were made (ASCII drawings for Haiku and
Fable, rasterised stroke geometry for Sonnet and Opus), in whether the templates
were centred and normalised the way MNIST is, and in the number of writing
styles per digit (one for Haiku, up to 47 slant and style variants for Opus).
Accuracy ranges from chance (Haiku) to 55.3\,\% (Opus). These are single samples
from each model; the earlier blind Fable network of Section~\ref{sec:mnist}
scored 55.4\,\% with the same approach, so run-to-run variance is of the order of
ten points and a model's benchmark score should be the mean of several runs
(\texttt{benchmark/RUNS.md}).

\begin{table}[h]\centering\small
\begin{tabular}{llrrrr}\toprule
task & version & labels seen & accuracy & params & non-zero \\\midrule
mnist & v1 blind & 0 & 55.4 & 5{,}890 & 544 \\
mnist & v2 rev1 & 500 & 74.3 & 5{,}890 & 553 \\
mnist & v3 rev2 & 500 & 79.4 & 7{,}066 & 662 \\
mnist & hole-finder & 5{,}000 & 83.3 & 21{,}772 & 2{,}411 \\
sst2 & v1 blind & 0 & 78.6 & 8{,}507 & 1{,}464 \\
sst2 & v2 rev1 & 872 & 79.2 & 9{,}167 & 1{,}567 \\
addition & v1 & 0 & 100.0 & 2{,}851 & 405 \\\midrule
\multicolumn{6}{l}{\emph{Cross-model run, standard prompt, blind, self-contained, one sample each}}\\
mnist & Claude Haiku 5.5 & 0 & 10.6 & 1{,}960 & 1{,}960 \\
mnist & Claude Sonnet 5.5 & 0 & 31.2 & 1{,}970 & 1{,}960 \\
mnist & Claude Fable 5.1 & 0 & 41.6 & 4{,}824 & 4{,}716 \\
mnist & Claude Opus 5.5 & 0 & 55.3 & 10{,}074 & 9{,}639 \\
\bottomrule\end{tabular}
\caption{HandNet-Bench rows (\texttt{benchmark/LEADERBOARD.md}). Top: this paper's networks as reference entries. Bottom: four models given the identical prompt. For the template MNIST networks the parameters are the drawings themselves.}
\label{tab:bench}
\end{table}
